% file: 1.函数与极坐标.tex
% 作为主文档 \input{1.函数与极坐标} 的子文件，不要写 \documentclass、\begin{document}、\end{document}
\usepackage{amsmath}
\usepackage{amssymb}

\section{函数的概念}
\subsection{实数与数轴}
实数的两个性质
\begin{enumerate}
    \item 有序性：任意两个互异的实数$a,b$都可比较大小，或者$a<b$，或者$a>b$。实数按照由小到大的顺序排列在数轴上。
    \item 完备性：因为任何两个有理点 $a,b$ 之间都有一个有理点 $\dfrac{a+b}{2}$，从而它们之间有无穷多个有理点，我们说有理点处处稠密。但有理点并未充满整个数轴，比如还有$\sqrt{2}$这样一些无理点。因为有理数与无理数之和为无理数，所以无理点也处处稠密。实际上，无理数远比有理数多得多。实数充满整个数轴，没有空隙，这就是实数的完备性(或连续性)。
\end{enumerate}

\noindent 数集与界

\noindent 绝对值

\noindent 函数的概念

\section{常用的概念}
\noindent 函数的特性

\noindent 隐函数和参数方程表示的函数

\noindent 单值函数与多值函数、反函数

\section{初等函数}
\subsection{基本初等函数及其图形}
\subsubsection{三角函数}
\[
\sin x=\frac{1}{\csc x}
\]
和差化积

积化和差

\subsection{复合函数与初等函数}
双曲正弦

双曲余弦

双曲正切

双曲余切

\section{极坐标}
\noindent 极坐标的概念

\noindent 极坐标与直角坐标的关系